\cotefile{91}
\scanpage{114}
% L1
\unsure{Posons|Pour}
% L2
\[
\operatorname{Cons}(X,E)\,(Y)
 = \operatorname{App}\,\operatorname{cons}(Y,(X,E)).
\]

% L3–4
\deleted{\unsure{On a les}}
 
Si \unsure{on a} $(X,E)$ et $(X',E')$ ; \unsure{alors on a}
% L5
\[
\operatorname{App}\,\operatorname{cons}((X,E),(X',E'))
 \xrightarrow{\sim}
 \deleted{\unsure{A}}\operatorname{Hom}(\operatorname{Cons}(X,E),\operatorname{Cons}(X',E'))
\]

% L6
On obtient ainsi \unsure{les|des} foncteurs pleinement fidèles
% L7–12
\diagram{\resizebox{\linewidth}{!}{\begin{tikzcd}[ampersand replacement=\&,column sep=large]
(\operatorname{Cons}\,\mathrm{aff})
 \arrow[r] \arrow[r,bend left=37,"{\substack{\text{\unsure{équivalence : on les|équivalence à des}}\\\text{\unsure{mor. surjectifs propres|ens. schématiques qu.-cpts}\ldots}}}"]
 \arrow[rr,bend right=39,"{\text{canonique}}"']
 \arrow[rrr,bend right=50,"{\text{can}}"']
\& \substack{(\text{Couples }(X,E))\\E\ \text{loc. const.}\\\text{sur }\deleted{X}\ \text{\unsure{quasicp|quasi-sép.}}} \arrow[r]
\& \substack{\operatorname{Ind}(\operatorname{Cons}\,\mathrm{aff})'\\
 \underbrace{\hphantom{\operatorname{Cons}\,\mathrm{aff}}}\\
 \text{systèmes}\\\text{inductifs}\\\text{\unsure{immersifs|inversifs|injectifs}}\ldots}
 \arrow[r] \arrow[r,bend right=43,"{\text{can}}"']
\& \underline{\operatorname{Hom}}((\operatorname{Cons}\,\mathrm{aff})^\circ,\mathrm{Ens})
\end{tikzcd}}}
% archive sheet number: 113
\scanpage{115}
% L1–3
Soit $\deleted{X}A$ un anneau commutatif, $X$ le schéma affine qu'il définit,
\[
 A^{\unsure{\mathrm{cons}|\mathrm{con}}}
 =\underset{\unsure{\mathscr{P}|P|\mathfrak P}}{\underline{\lim}}\ A_{\mathscr{P}}
\]
% L4–6
$\mathscr{P}$ parcourt les ``partitions \unsure{standard|standart}'' de $X$ en ens constructibles loc. fermés \unsure{,|et} affines, et pour une telle $\mathscr{P}=(X_i)_{i\in I}$, $A_{\mathscr{P}}$ désigne \deleted{\unsure{le produit}} $\unsure{A|\Gamma}(\coprod X_{i\,\mathrm{red}})=\deleted{\unsure{\mathscr{P}}}\prod A(X_{i\,\mathrm{red}})$

% L7–8
Les hom de transition sont injectifs. Considérons
\[
 \widetilde X=\operatorname{Spec}(A^{\mathrm{cons}})
\]
% L9
On a $\operatorname{Top}(\widetilde X)=\underset{\mathscr P}{\underleftarrow{\operatorname{Lim}}}\operatorname{Top}\operatorname{Spec}(A_{\mathscr P})$

% L10–11
En particulier
\[
 \operatorname{Ens}(\widetilde X)=\underset{\mathscr P}{\underleftarrow{\operatorname{Lim}}}\ 
 \underbrace{\operatorname{Ens}\operatorname{Spec} A_{\mathscr P}}_{\coprod X_i}
\]
% L12–14
et comme les morphismes de transition sont des \deleted{\unsure{homéomorphismes}}\inserted{bijections}, on trouve \deleted{que} une \deleted{\unsure{surj}}\inserted{bijection}.
\[
 \operatorname{Ens}\widetilde X\simeq\operatorname{Ens} X
\]
% L15–18
Cette bijection est compatible avec les topologies de $\widetilde X$ et celle de $X^{\mathrm{cons}}$, comme --- on voit aisément sur la description d'une base d'ouverts de $\widetilde X$ \ldots
% archive sheet number: 114
\scanpage{116}
% L1–2; large diagonal strokes through the entire draft
\deleted{Soit $X$ un espace topologique, \deleted{\unsure{dont chaque pt a un syst.|dont chaque pt $x$ a un syst.}}\inserted{ayant une base d'ouverts} \deleted{fondamental de voisinages} $U$ ayant les propriétés suivantes :
\margin{\unsure{Hyp|1\textsuperscript{ère}}}

% L3
(i) $U$ rétrocompact dans $X$, et $U$ quasi-compact

% L4–6
(ii) T\textsuperscript{te} \deleted{suite} famille $(E_i)_{i\in I}$ de parties \deleted{\unsure{fermées}} constructibles de $U$, recouvrant $U$, il y a une sous-famille finie recouvrant $U$.

% L7–11
Pour exprimer (ii), Introduisons [sur l'espace topologique $U$] la topologie dont les ouverts sont les \inserted{réunions quelconques de} parties loc. constructibles [c'est bien stable par intersections finies et réunions \unsure{qlc|qlq|quelc.}] donc une base d'ouverts formée des parties constructibles.
% L12–17
Lorsque $U$ est quasi-compact, et \deleted{que} \deleted{\unsure{$U$ qui admet|$U$ est admet}} admet une base d'ouverts rétrocompacts, ``loc constructible'' est d'ailleurs pareil que ``constructible'' ; Moyennant (i), les $U$ intervenant dans (Hyp) tombent dans le cas précédent.
\margin{$\varnothing$.}
% L17–19
Notons \inserted{que} si de plus $\deleted{\unsure{X}}U$ satisfait l'axiome $T_0$, \inserted{[c'est le cas des préschémas \unsure{\ldots}]} \deleted{\unsure{$U$ est}} alors une telle topologie est \uline{séparée}, et \unsure{même|en} totalement discontinue [pour deux pts $x\ne y$, il existe des ouverts $V\ni x$, $W\ni y$, $V\cap W=\varnothing$, $V\cup W=U$]
% L20–22
L'hyp (ii) exprime qu'elle est \uline{quasi-compacte} [donc \uline{compacte} dans le cas particulier que nous envisageons].}
% archive sheet number: 115
\scanpage{118}
% L1–2
\uline{N\textsuperscript{o} 1} Préschémas \unsure{cpts|cpcts} et loc. cpts. Anneaux flasques.

\uline{N\textsuperscript{o} 2} \deleted{Systèmes projectifs}\inserted{Proschémas} flasques \deleted{\unsure{de préschémas et} \unsure{constructibisation} \illegible{3}}

% L3–4
\uline{N\textsuperscript{o} 3} \deleted{Espaces \deleted{\unsure{ind}} schématiques Ind = \unsure{Espaces|ensembles}}\inserted{espaces schématiques compacts} Espaces schématiques. [Catégorie \unsure{semblable pour|véritable pour|engendrée par}] les rec. ouverts \ldots{} ]
\margin{Morphismes quasi-compacts $\rightarrow$ \unsure{séparation|opérations} propres}
% slanted

% L5
\uline{N\textsuperscript{o} 4} Opérations sommes, \deleted{\unsure{sous}}\inserted{\unsure{cas des quasi-compacts|des proschémas} $\simeq$ sommes de compacts} et $\underleftarrow{\lim}$ \deleted{\unsure{finies}} des espaces schématiques

% L6
\uline{N\textsuperscript{o} 5} \deleted{\unsure{Esp. quasi-schématiques}} Espaces quasi-schématiques \ldots\ldots

% L7–9: cancelled diagram
\deleted{\diagram{\resizebox{\linewidth}{!}{\begin{tikzcd}[ampersand replacement=\&,column sep=large]
\left(\substack{\text{Préschémas quasi-compacts}\\\text{quasi-séparés}}\right)\arrow[r,"\alpha"]
\&\left(\substack{\text{Esp. schématiques}\\\text{compacts}}\right)
 \arrow[r,"{\substack{\inserted{\text{équiv.}}\\\deleted{\unsure{\text{fidèle}}}\\\approx}}"]
\&\left(\operatorname{Pro}(\mathrm{Sch})\ \substack{\text{flasques}\\\deleted{\unsure{\text{finis}}}\ \text{q.cpts}}\right)
 \arrow[r,"{\substack{\text{fid, }\inserted{\text{submersif}}\\\deleted{\unsure{\text{ess. surjectif}}}}}"]
\&\left(\substack{\text{schémas compacts}\\\text{réduits}}\right)
 \arrow[l,bend left=24,"{\text{pl. fid. }\unsure{\sigma|\gamma}}"']
 \arrow[lll,bend left=32,"{\text{pl. fid.}}"]
\end{tikzcd}}}}

% L10
\uline{N\textsuperscript{o} 6} \unsure{pts|pbs|Pb} épais. \unsure{Espaces|Espace} de Néron. \deleted{\illegible{1}}

% L11–15
\diagram{\resizebox{\linewidth}{!}{\begin{tikzcd}[ampersand replacement=\&,column sep=huge]
\substack{(\text{Préschémas quasi-compacts}\\\text{quasi-séparés})\\[3pt] X}
 \arrow[r,"{\alpha}","{\substack{\text{submersif}\\\text{[fidèle sur les}\\\text{réduits].}}}"']
 \arrow[rr,bend left=24,"{\sigma}"]
\&\substack{\text{(Cons cpt)}\\ \shortparallel\\(\operatorname{Pro}(\mathrm{Sch})\ \text{flasques})\\[3pt]\mathscr Y}
 \arrow[r,"{\beta=\underleftarrow{\lim}}","{\substack{\text{submersif}\\\text{fidèle}}}"']
\&\substack{(\text{schémas compacts}\\\text{réduits})\\[3pt]K}
 \arrow[l,bend left=24,"{\beta'\ \text{pl. fid}}"']
 \arrow[ll,bend left=33,"{\sigma'\ \text{pl. fid. }\unsure{\text{imm. can.}|\text{imm. ouv.}|\text{inv. can.}|\text{immersion}}}"']
\end{tikzcd}}}
\margin{$\nearrow$ [pleinement fidèle également sur les ``de type fini'' \unsure{en un sens|au sens}\ \unsure{précisé\ldots|fonc\ldots}]}
% L15–16
$\beta'(K)={}$ système proj trivial défini par $K$

% L17–18
\[
\begin{aligned}
\operatorname{Hom}(K,\sigma X)&\xrightarrow{\sim}\operatorname{Hom}(\sigma'K,X)
&\operatorname{Hom}(K,\beta(\alpha X))&=\operatorname{Hom}(\beta'K,\alpha X)\\
\operatorname{Hom}(K,\beta\mathscr Y)&\xrightarrow{\sim}\operatorname{Hom}(\beta'K,\mathscr Y)
\end{aligned}
\]
% archive sheet number: 117
\scanpage{119}
\mspage{- 90 -}
n\textsuperscript{o} 371

% L1–2
tout $i$. Montrons que l'on a $(u(x))^2=0$ pour tout $x\in M$ ; on peut en effet écrire $x=\sum_{i=1}^{n}x_i$, où $x_i\in M_i$, donc $(u(x))^2=\sum_{i=1}^{n}(u(x_i))^2+\sum_{i<j}(u(x_i)u(x_j)+u(x_j)u(x_i))$.
% L3–6
Or, comme $\handwritten{\psi}_i(x_i)$ est un élément de degré 1 dans $\handwritten{\bigwedge}(M_i)$, on a, pour $i\ne j$, $h_i(\handwritten{\psi}_i(x_i))h_j(\handwritten{\psi}_j(x_j))=-h_j(\handwritten{\psi}_j(x_j))h_i(\handwritten{\psi}_i(x_i))$ (\S~3, n\textsuperscript{o}3, prop.2), d'où $u(x_i)u(x_j)+u(x_j)u(x_i)=0$ ; par ailleurs $(u(x_i))^2=h_i((\handwritten{\psi}_i(x_i))^2)=0$ par définition de $\handwritten{\bigwedge}(M_i)$. En vertu de la prop.3 du n\textsuperscript{o}2, il existe donc un homomorphisme unique d'algèbres graduées
% L7–8
$h:\handwritten{\bigwedge}(M)\longrightarrow {}^{\mathrm g}\!\bigotimes_i(\handwritten{\bigwedge}(M_i))$ tel que $u=h\circ\handwritten{\psi}$ (où $\handwritten{\psi}:M\longrightarrow\handwritten{\bigwedge}(M)$ est l'injection canonique). Pour $x_i\in M_i$, on a
% L9–12
$h(g(u_i(x_i)))=h(g(h_i(\handwritten{\psi}_i(x_i))))=h(f_i(\handwritten{\psi}_i(x_i)))=h(\handwritten{\psi}(j_i(x_i))=u_i(x_i)$, et comme les $u_i(x_i)$ engendrent l'algèbre $ {}^{\mathrm g}\!\bigotimes_i(\handwritten{\bigwedge}(M_i))$, $h\circ g$ est l'identité ; de même $g(h(\handwritten{\psi}(j_i(x_i))))=g(u_i(x_i))=g(h_i(\handwritten{\psi}_i(x_i)))=f_i(\handwritten{\psi}_i(x_i))=\handwritten{\psi}(j_i(x_i))$, et les éléments $\handwritten{\psi}(j_i(x_i))$ engendrent l'algèbre $\handwritten{\bigwedge}(M)$, donc $g\circ h$ est l'identité.

% L13–16
COROLLAIRE 1.- \uline{Sous les hypothèses de la prop.7}, \uline{pour tout entier} $m\geq 0$, \uline{la puissance extérieure m-ème} $\overset{m}{\handwritten{\bigwedge}}\,M$ \uline{s'identifie canoniquement par} $g$ \uline{à la somme directe des modules} $\bigotimes_{1\leq i\leq n}(\overset{m_i}{\handwritten{\bigwedge}}\,M_i)$, \uline{où} $(m_i)_{1\leq i\leq n}$ \uline{parcourt l'ensemble des suites d'entiers} $\geq 0$ \uline{telles que} $\sum_{i=1}^{n}m_i=m$.

% L17–19
Cet isomorphisme s'explicite de la façon suivante : considérons $m$ éléments $x_j=\sum_{i=1}^{n}x_{ji}$ $(1\leq j\leq m)$ de $M$, avec $x_{ji}\in M_i$ pour $1\leq i\leq n$, et pour toute application $\handwritten{\sigma}$ de $\unsure{[1,m]|[\![1,m]\!]}$ dans $[1,n]$, formons le m-vecteur
% L20
\begingroup\csname tagsleft@true\endcsname
\[
 x_{1,\handwritten{\sigma}(1)}\handwritten{\wedge} x_{2,\handwritten{\sigma}(2)}\handwritten{\wedge}\ldots\handwritten{\wedge} x_{m,\handwritten{\sigma}(m)}\quad .\tag*{(7)}
\]
\endgroup
% L21–24
La donnée de l'application $\handwritten{\sigma}$ équivaut à la donnée, pour chaque $i\in[1,n]$ d'une partie $H_i=\handwritten{\sigma}^{-1}(i)$ de $[1,m]$, celles des $H_i$ non vides formant une \uline{partition} de $[1,m]$ ; si $H_i$ est vide, posons $x_{H_i,\,i}=1$, et dans le cas contraire, $x_{H_i,\,i}=x_{j_1,\,i}\handwritten{\wedge} x_{j_2,\,i}\handwritten{\wedge}\ldots\handwritten{\wedge} x_{j_{r_i},\,i}$, où $(j_k)_{1\leq k\leq r_i}$ est la suite
% archive sheet number: 118
\scanpage{120}
% L1–3
\diagram{\resizebox{\linewidth}{!}{\begin{tikzcd}[ampersand replacement=\&,column sep=large]
(\text{Préschémas quasi-séparés}) \arrow[r,"\bar\alpha"]
 \arrow[rr,bend left=14,"{\bar{\unsure{\sigma|\gamma}}}"]
\&\substack{\operatorname{Ind}(\operatorname{Cons}\,\mathrm{cpt})\ \text{flasques}\\\shortparallel\\(\operatorname{Cons})}
 \arrow[r,"{\bar\beta}","{\text{submersif}}"']
\&\left(\substack{\text{schémas loc. cpts}\\\text{réduits}}\right)
 \arrow[l,bend left=21,"{\bar\beta'\ \text{pl fid}}"']
 \arrow[ll,bend left=28,"{\bar\sigma'\ \text{pl. fid}}"']
\end{tikzcd}}}

% L4–7: the thin line runs from the second Hom to the second row of limits.
\diagram{\resizebox{\linewidth}{!}{\begin{tikzcd}[ampersand replacement=\&,column sep=small,row sep=large]
\operatorname{Hom}(\bar K,\bar\beta\,\overline{\mathscr Y})
 \simeq\operatorname{Hom}(\bar\beta'(K),\overline{\mathscr Y})
 \arrow[dr,dash,out=-8,in=170]
\&{}\simeq\underset{\lambda}{\underleftarrow{\lim}}\ \underset{\mu}{\underrightarrow{\lim}}
 \operatorname{Hom}(\bar K_\lambda,\beta(\overline{\mathscr Y}_\mu))\\
\deleted{\operatorname{Hom}(\bar K,\bar\sigma(X))\simeq}
\&{}\simeq\underset{\lambda}{\underleftarrow{\lim}}\ \underset{\mu}{\underrightarrow{\lim}}
 \operatorname{Hom}(\beta'(\bar K_\lambda),\overline{\mathscr Y}_\mu)\\
\operatorname{Hom}(\bar K,\bar\sigma(X))\simeq\operatorname{Hom}(\bar\sigma'K,X)
\&{}\simeq\underset{\lambda}{\underleftarrow{\lim}}\ \underset{\mu}{\underrightarrow{\lim}}
 \operatorname{Hom}(K_\lambda,\sigma(X_\mu))\\
\&{}\simeq\underset{\lambda}{\underleftarrow{\lim}}\ \underset{\mu}{\underrightarrow{\lim}}
 \operatorname{Hom}(\sigma'K_\lambda,\deleted{\unsure{\sigma}}X_\mu)
\end{tikzcd}}}

% L8
\[
 (\operatorname{Cons})\xrightarrow{\text{pl. fid}}\underrightarrow{\operatorname{Hom}}((\operatorname{Cons}\,\mathrm{cpt})^\circ,(\operatorname{Cons}))
\]
% archive sheet number: 119
\scanpage{121}
\mspage{- 93 -}
n\textsuperscript{o} 371

% L1–2
$\handwritten{\bigwedge}(N)$, où $N$ est un sous-module de type fini de $\unsure{M}$ ; il suffit alors d'appliquer le cor.3 de la prop.2 du n\textsuperscript{o}1.

% L3–6
COROLLAIRE 5.- \uline{Soient} $N$, $P$ \uline{deux sous-modules d'un} $A$-\uline{module} $M$ \uline{tels que} $N\cap P$ \uline{soit facteur direct de} $N$ \uline{et de} $P$ \uline{et} $N+P$ \uline{facteur direct de} $M$. \uline{Alors} $\handwritten{\bigwedge}(N)$ \uline{et} $\handwritten{\bigwedge}(P)$ \uline{s'identifient canoniquement à des sous-algèbres de} $\handwritten{\bigwedge}(M)$, \uline{et} $\handwritten{\bigwedge}(N\cap P)$ \uline{à l'intersection} $\handwritten{\bigwedge}(N)\cap \handwritten{\bigwedge}(P)$.

% L7–12
En effet, si $E$, $F$ sont des supplémentaires de $N\cap P$ dans $N$ et $P$ respectivement et $G$ un supplémentaire de $N+P$ dans $M$, $M$ est somme directe des quatre sous-modules $N\cap P$, $E$, $F$, $G$ ; donc $\handwritten{\bigwedge}(N)$, $\handwritten{\bigwedge}(P)$ et $\handwritten{\bigwedge}(N\cap P)$ s'identifient à des sous-algèbres de $\handwritten{\bigwedge}(M)$ (cor.3). En outre, en vertu du cor.1, $\overset{n}{\handwritten{\bigwedge}}\,N$ (resp. $\overset{n}{\handwritten{\bigwedge}}\,P$) s'identifie à la somme directe des sous-modules $\overset{p}{\handwritten{\bigwedge}}(N\cap P)\otimes\overset{n-p}{\handwritten{\bigwedge}}\,E$ (resp. $\overset{p}{\handwritten{\bigwedge}}(N\cap P)\otimes\overset{n-p}{\handwritten{\bigwedge}}\,F$) pour $0\leq p\leq n$ ; on a donc bien $\handwritten{\bigwedge}(N\cap P)=\handwritten{\bigwedge}(N)\cap \handwritten{\bigwedge}(P)$.

% L13–14
PROPOSITION 8.\unsure{=|-} \uline{Soit} $(e_{\lambda})_{\lambda\in L}$ \uline{une base d'un} $A$-\uline{module} $M$ ; \uline{munissons} $L$ \uline{d'une structure d'ordre total}, \uline{et pour toute partie finie} $H$ \uline{de} $L$, \uline{posons}
% L15
\begingroup\csname tagsleft@true\endcsname
\[
 e_H=e_{\lambda_1}\handwritten{\wedge} e_{\lambda_2}\handwritten{\wedge}\mathord{..}\handwritten{\wedge} e_{\lambda_n}\tag*{(11)}
\]
\endgroup
% L16–18
\uline{si} $(\lambda_i)_{1\leq i\leq n}$ \uline{est la suite strictement croissante des éléments de} $H$. \uline{Alors les éléments homogènes} $e_H$, \uline{où} $H$ \uline{parcourt l'ensemble} $F(L)$ \uline{des parties finies de} $L$, \uline{forment une base de} $\handwritten{\bigwedge}(M)$.

% L19–23
Supposons d'abord $L$ \uline{fini}, soit par exemple $L=\unsure{[1,m]|[\![1,m]\!]}$ muni de l'ordre usuel ; alors la proposition résulte de la prop.7, car si on pose $M_i=Ae_i$ pour $1\leq i\leq m$, $M$ est somme directe des $M_i$ et il résulte aussitôt du cor.3 de la prop.2 du n\textsuperscript{o}1 que $\handwritten{\bigwedge}(M_i)=A\oplus M_i$ (autrement dit $\overset{p}{\handwritten{\bigwedge}}M_i=0$ pour $p\geq 2$) ; la conclusion résulte alors du cor.1 de la prop.7 et du chap.II, 3e éd., \S~3, n\textsuperscript{o}4, prop.4.

% L24–25
Supposons ensuite $L$ infini et pour toute partie finie $J$ de $L$, soit $M_J$ le sous-module de $M$ ayant pour base les $e_\lambda$ tels que $\lambda\in J$ ; il est clair que $M$ est limite
% archive sheet number: 120
\scanpage{122}
% L1
$X$ Schéma compact, [i.e. quasi-compact et \inserted{\uline{topologiquement}} \uuline{séparé}]

% L2–6
(i) $X$ est tot. discontinu, \deleted{i.e. si $x,y\in X$, $x\ne y$, alors il existe une partie ouverte et fermée de $X$ qui contient $x$ et non $y$. En effet, soient $U,V$ des ouverts, $U\ni x$, $V\ni y$, $U\cap V=\varnothing$.} i.e. tt $x\in X$ a un syst. fond. de voisinages qui sont à la fois \uline{ouverts} et \uline{fermés}

% L7–8
[En effet, il y a un syst. fond. de voisinages ouverts quasi compacts]

% L9–11
(ii) $X$ \unsure{est|\deleted{est}} affine. En effet, \deleted{\unsure{on construit de proche en proche des ouverts affines non vides qui ne rencontrent pas}}\inserted{on recouvre $X$ par des ouverts affines en nb fini, soit $X_i$ $(1\leq i\leq n)$}. Ils sont ouverts et fermés. Soit
% L12–16
$X'_i=X_i-X_i\cap\bigcup_{j<i}X_j$, Alors $X'_i$ est une partie ouverte et fermée de $X_i$, donc affine puisque $X_i$ l'\unsure{est} (en vertu p. ex. du critère \unsure{coh.|ult.} \deleted{\unsure{,}} ou \unsure{n'imp. quoi|n'imp. quel}). \unsure{Comme|Donc} $X$ est \deleted{\unsure{est}} la somme des $X'_i$, il est affine.

% L17–19
(iii) Les parties constructibles de $X$ sont ses parties ouvertes et fermées. [\deleted{Donc}\inserted{\unsure{En dts|En ft|En effet}} $\operatorname{Cons} X\to X$ est un \deleted{\unsure{homéom.}]} \unsure{isom|iso.|bon}] Donc
% archive sheet number: 121
\scanpage{123}
\mspage{- 92 -}
n\textsuperscript{o} 371

% L1–2
\uline{de sous-modules} ; \uline{pour tout} $i$, \uline{soit} $u_i=M_i\to N_i$ \uline{une application} $A$-\uline{linéaire et soit} $u=\bigoplus_{1\leq i\leq n}u_i:M\longrightarrow N$. \uline{Alors le diagramme}
% L3–5
\diagram{\begin{tikzcd}[ampersand replacement=\&,column sep=large,row sep=large]
{}^{\mathrm g}\!\bigotimes_{1\leq i\leq n}(\handwritten{\bigwedge}(M_i)) \arrow[r] \arrow[d,"{{}^{\mathrm g}\!\bigotimes_{1\leq i\leq n}(\handwritten{\bigwedge}(u_i))}"']
 \& \handwritten{\bigwedge}(M)\arrow[d,"{\handwritten{\bigwedge}(u)}"]\\
{}^{\mathrm g}\!\bigotimes_{1\leq i\leq n}(\handwritten{\bigwedge}(N_i))\arrow[r] \& \handwritten{\bigwedge}(N)
\end{tikzcd}}
% L6–7
(\uline{où les flèches horizontales sont les isomorphismes définis dans la prop.7}) \uline{est commutatif}.

% L8–10
Il suffit de voir que les deux applications $ {}^{\mathrm g}\!\bigotimes_i(\handwritten{\bigwedge}(M_i))\longrightarrow \handwritten{\bigwedge}(N)$ définies par le diagramme prennent la même valeur dans chacun des modules $\bigotimes_{1\leq i\leq n}(\overset{m_i}{\handwritten{\bigwedge}}M_i)$, ce qui résulte aussitôt des \unsure{définitions|défintions}.

% L11–14
COROLLAIRE 3.- \uline{Soit} $u:M\to N$ \uline{une application} $A$-\uline{linéaire injective}, \uline{telle que} $u(M)$ \uline{soit facteur direct de} $N$. \uline{Alors l'homomorphisme} $\handwritten{\bigwedge}(u):\handwritten{\bigwedge}(M)\longrightarrow \handwritten{\bigwedge}(N)$ \uline{est injectif et} $\handwritten{\bigwedge}(N)$ \uline{est somme directe de la sous-algèbre graduée} $\handwritten{\bigwedge}(u)(\handwritten{\bigwedge}(M))$ \uline{et d'un idéal bilatère gradué engendré par un supplémentaire de} $u(M)$ \uline{dans} $N$.

% L15–20
Cela résulte du cor.1 appliqué pour $n=2$, car la somme des sous-modules $(\overset{p}{\handwritten{\bigwedge}}\,M_1)\otimes(\overset{q}{\handwritten{\bigwedge}}\,M_2)$ pour tous les couples $(p,q)$ tels que $p\geq 1$ n'est autre que l'idéal bilatère $m_1\otimes(\handwritten{\bigwedge}(M_2))$ de $(\handwritten{\bigwedge}(M_1))\otimes^{\mathrm g}(\handwritten{\bigwedge}(M_2))$, où $m_1$ est la somme directe des $\handwritten{\bigwedge}(M_1)$ pour $p\geq 1$, idéal bilatère de $\handwritten{\bigwedge}(M_1)$. On peut aussi donner une démonstration directe n'utilisant pas la prop.7 et calquée sur celle du \S~4, n\textsuperscript{o}2, prop.4.

% L21–22
COROLLAIRE 4.- \uline{Tout élément} $z$ \uline{de} $\handwritten{\bigwedge}(M)$ \uline{dont la composante homogène scalaire} (i.e. \uline{de degré} 0) \uline{est nulle}, \uline{est nilpotent}.

% L23
Il est clair en effet que tout élément de $\handwritten{\bigwedge}(M)$ appartient à une sous-algèbre
% archive sheet number: 122
\scanpage{124}
% L1
\hfill\uline{Greenberg} - \uline{Néron}
% archive sheet number: 123
\scanpage{126}
% L1–3
Soit $\mathscr X=(X_i)_{i\in I}$ un système projectif de préschémas, essentiellement affine [i.e. $\exists\,i_0,\ \forall\,i\geq i_0,\ \unsure{\varphi|\phi}_{i_0i}$ est affine]. Alors $X=\underset{i}{\underleftarrow{\lim}} X_i$ existe, $\operatorname{Top} X=\underset{i}{\underleftarrow{\lim}}\operatorname{Top}(X_i)$ \ldots.

% L4
Soit $x\in X$, on dit que $\unsure{X|x}$ est \uline{épais} si
% L5–7
\begingroup\csname tagsleft@true\endcsname
\[
\begin{aligned}
&\exists\,i\quad\text{t.q.}\quad\forall\,j\geq i,\quad x_j\text{ est un pt maximal}\\
&\text{de }\varphi_{\deleted{\unsure{j}}ij}^{-1}(x_i),\quad\text{donc si }k\geq j\geq i,\quad\deleted{\unsure{\varphi_{ij}}}x_k\text{ est}\\
&\text{maximal de }\varphi_{jk}^{-1}(x_j).
\end{aligned}\tag*{(Ep)}
\]

% L8
Sous quelles conditions cela équivaut :
% L9–11
\[
\begin{aligned}
&\exists\,i,\quad\deleted{\text{tq}}\ \text{et une partie ouverte non vide }U_i\text{ de }\bar x_i,\\
&\text{t.q.}\quad\forall\,j\geq i,\quad\varphi_{ij}^{-1}(U_i)\subset\bar x_j\qquad\text{(alors la}\\
&{}=\text{ propriété est vraie pour tt }i'\geq i\text{).}
\end{aligned}\tag*{(Ép\textsuperscript{\unsure{?|1}})$'$}
\]

\endgroup
% L12–13
Cela prouve que $\varphi_{ij}^{-1}(U_i)$ est irréductible, de pt générique $x_j$, et donc $(\mathrm{Ép})'$ implique $(\mathrm{Ep})$\unsure{?|!}

% L14
Considérons aussi
% L15–17
$(\mathrm{Ep})''$\quad $\bar x\inserted{{}=\underleftarrow{\lim}\bar x_i}$ est une partie \uline{constructible} de $X$.

Ceci équivaut \inserted{(si $X_i$ sont \unsure{noeths|noeth.})} à : $\bar x=\varphi_i^{-1}(\bar x_i)$ pour $i$ grand \deleted{\unsure{(ou : pour $i$ convenable \ldots)|(Cons. : pour $i$ constructible \ldots)}} (donc pour $i'\geq i$)

% L18–19
$(\mathrm{Ep})'''$\quad $\exists$ une partie constructible \deleted{loc. fermée}\inserted{ouverte non vide} de $\bar x$, soit $U$.

i.e. $\exists\,U_i$
% archive sheet number: 125
\scanpage{127}
\mspage{- 145 -}
n\textsuperscript{o} 316

% L1–3
Si $x\longrightarrow x_M$ est une représentation irréductible de dimension finie, si $\handwritten{\beta}$ est la forme bilinéaire associée, et si $\dim h$ n'est pas divisible par la caractéristique de $K$, montrer que $H^p(g,M)=\{0\}$
% L4–7
pour tout $p$ . (Utilisant la prop.12 , montrer que $\unsure{\Gamma|T}$ est un automorphisme de $C^{*}(g,M)$ , donc induit des automorphismes de $Z^p(g,M)$ et $B^p(g,M)$ . Pour $f\in Z^p(g,M)$, $\Gamma f={}$ $=(d\handwritten{\unsure{\varrho|\rho}}+\handwritten{\varrho}d)f=d\handwritten{\varrho}f\in B^p(g,M)$ , d'où $Z^p(g,M)=B^p(g,M)$).
\par\centerline{\rule{3cm}{0.4pt}}
% archive sheet number: 126
\scanpage{128}
% L1: above the frame
$x_i$ maximal dans $\varphi_{\unsure{ji|ij}}^{-1}(\deleted{\unsure{x}}x_{\unsure{i|j}})$

% L2–5: implication diagram, with the lower left node outside the outline
\diagram{\resizebox{\linewidth}{!}{\begin{tikzcd}[ampersand replacement=\&,column sep=2.4em,row sep=1.6em,execute at end picture={\draw[rounded corners] ($(epone.north west)+(-.35,.25)$) rectangle ($(epthree.south east)+(4.65,-.24)$);}]
\& |[alias=epone]|(\mathrm{Ep})' \arrow[rr,Rightarrow] \arrow[dr,Rightarrow] \& \& (\mathrm{Ep})+\bigl(\underline{\text{unicité : \unsure{parler|problème|probl.} de }x_i}\bigr)\\
\& \deleted{(\mathrm{Ep})'''} \& |[alias=epthree]|(\mathrm{Ep})''' \arrow[ur,Rightarrow,"{\substack{\text{si }\deleted{\unsure{\forall}}\,\forall x_i\in\deleted{\unsure{X|X_i}}X'_i,\ \exists\text{ un ouvert }\ne\varnothing\\ U_i\subset\bar{x}_i\text{ t.q. les }\varphi_{ij}^{-1}(U_i)\cap X'_j=\varphi_{\unsure{ij|j}}(X_i)}}"{description,at end,anchor=north west}] \&\\
(\mathrm{Ep})'' \arrow[uur,Rightarrow,"{\substack{\text{M.L.}\\\text{et }I=\underline{N}}}"] \arrow[urr,Rightarrow,"{4}"'] \& \& \&
\end{tikzcd}}}

\[\text{???}\]

% L2–4: right
$\deleted{\unsure{E|P}}$\qquad
\diagram{\begin{tikzcd}[ampersand replacement=\&]
\text{pt épais} \& \substack{\text{parties constructibles}\\\text{irréductibles}} \arrow[l]
\end{tikzcd}}
\[
\begin{aligned}
x&\in E\subset\bar{x}\\
x&\in U_0\cap Z_0\\
&\bar{x}
\end{aligned}
\]

% L6–8: left and centre
\[
\begin{aligned}
U_j\cap X'_j&\subset\bar{x}_j\\
U_i&\subset\bar{x}_i
\end{aligned}
\]
\diagram{\begin{tikzcd}[ampersand replacement=\&]X'_i\&X'_j\arrow[l]\end{tikzcd}}
\[
x\in\deleted{\unsure{Z_x|Z_0|\bar{x}}}\,\bigcup_i U_i\cap Z_i
\qquad x\in U_0\cap Z_0\subset\bar{x}
\]
% L9–10
\[
\deleted{\bar{x}\in}\quad\deleted{x\in Z_{\unsure{0|{}}}}
\qquad \bar{x}=\overline{U}_0\cap\unsure{\widetilde{Z}_0|Z_0}\subset Z_0
\]
\[\deleted{U_0\subset}\]
\diagram{\begin{tikzcd}[ampersand replacement=\&,column sep=0.4em]
\mathscr X\arrow[dr]\&\&\mathscr X_i\arrow[dl]\\
\&{}\&
\end{tikzcd}}

% L11–15: lower right
$f_1\quad f_n$ nuls en $x$
\[
f_1 \mathinner{\ldotp\ldotp\ldotp} f_n\in\mathfrak p\qquad g
\]
\[
\text{si } g\notin\mathfrak q \qquad \mathfrak q\supset\mathfrak p \Longleftrightarrow
\mathfrak q\supset f_1\mathinner{\ldotp\ldotp} f_n
\]
\[
\text{\textcircled{1}}\qquad \underline{O}_{X,x}\qquad
\mathfrak m_x\quad\deleted{\unsure{\text{engendré}}}
\]
$A_g$ a un \unsure{idéal} de \unsure{définition}\qquad
\unsure{idéal de} \deleted{\unsure{type}} \unsure{par} un \unsure{nb} fini d'\unsure{élts}

% archive sheet number: 127
